\section{Literature Review}
\label{ch:literature}

This section synthesises the core literature that motivates the present
computational study. The review is organised around project logic rather than
paper-by-paper summary: working memory as recurrent dynamics; trained RNNs as
testbeds; acute psilocybin behavioural signatures; psychedelic neural-dynamics
theories; and serotonergic gain mechanisms. Direct evidence, modelling
analogies, and speculative translations are kept distinct.

\subsection{Working memory as recurrent dynamics}

Working memory can be supported by recurrent excitation that sustains
task-relevant activity after sensory drive has disappeared. In prefrontal
cortex, NMDA-dependent recurrent mechanisms contribute to delay-period
persistent firing \cite{wang2013nmda}. At the population level, however,
stable coding need not require every neuron to fire tonically: Murray and
colleagues showed that heterogeneous single-unit trajectories can coexist with
a low-dimensional subspace in which stimulus representations remain decodable
across cue and delay epochs \cite{murray2017stable}. Theoretical work likewise
emphasises that recurrent heterogeneity and attractor structure shape drift,
noise robustness, and memory stability \cite{kilpatrick2013}.

Trained recurrent networks provide a complementary modelling route. Song,
Yang, and Wang showed that excitatory--inhibitory RNNs can be trained on
cognitive tasks, including parametric working memory, while remaining useful
as mechanistic hypothesis generators provided that dynamics and connectivity
are analysed rather than behaviour alone \cite{song2016ei}. Ghazizadeh and
Ching further showed that trained working-memory networks can encode
information along slow manifolds, with multiple dynamical solutions capable of
supporting delay-period performance and different robustness/forgetting
tradeoffs \cite{ghazizadeh2021}. Frontoparietal circuit models also motivate
testing maintenance under distraction and distributed recurrent roles
\cite{murray2017frontoparietal}. Together, these sources justify a
task-trained RNN baseline judged by both behavioural competence and
population-level delay dynamics.

\subsection{Acute psilocybin effects on attention and working memory}

Human psilocybin cognition does not reduce to a single global impairment.
Across attention and executive tasks, a multilevel meta-analysis found a large
acute slowing of reaction time (Hedges' $g=1.13$, 95\% CI $[0.57,1.70]$) with
a weaker and less robust accuracy effect \cite{yousefi2025}. Dose dependence
and task sensitivity were important: broader or more demanding measures tended
to show clearer slowing. This is the strongest behavioural rationale for
including timing-sensitive outcomes, such as response settling, alongside final
accuracy.

Original studies refine the selective pattern. In a letter N-back task,
psilocybin impaired updating more clearly at 2-back than at 0-back, with
reduced discriminability and prolonged responding \cite{barrett2018}. In a
small within-subject study, multiple-object tracking was impaired while spatial
span was comparatively spared \cite{carter2005}. Interval reproduction and
spatial span showed stronger impairment at higher dose and longer intervals
\cite{wittmann2007}. On a memory-guided delayed-response task, reaction time
increased without a clear reduction in correct responses
\cite{vollenweider1998}. Broader reviews likewise report mixed working-memory
findings and more consistent acute attentional vulnerability, while cautioning
against equating acute healthy-volunteer effects with post-acute clinical
outcomes \cite{meshkat2024,sayalibarrett2023}. Sayali and Barrett additionally
frame psychedelic cognition as a stability--flexibility tradeoff rather than
uniform deficit \cite{sayalibarrett2023}.

For modelling, these findings support a selective target: greater cost under
updating load and distraction, response slowing with relative preservation of
simple accuracy, and stronger impairment at greater perturbation strength.
They do not specify an RNN equation.

\subsection{Psychedelic neural-dynamics theories and models}

Theoretical accounts motivate altered precision, gain, and constraint. The
entropic-brain framework emphasised increased signal diversity and reduced
ordered constraint under psychedelics \cite{carhartharris2014}. REBUS relates
relaxed high-level prior precision to altered postsynaptic gain and greater
influence of afferent prediction-error signals \cite{carhartharris2019}. These
are interpretive frameworks, not working-memory implementations.

Computational and empirical models sharpen the translation space. Herzog and
colleagues operationalised 5-HT2A effects in a whole-brain dynamic mean-field
model as receptor-density-dependent modulation of the neural response function,
producing heterogeneous entropy increases shaped by recurrent topology
\cite{herzog2023}. Empirical MEG work similarly reports increased spontaneous
signal diversity under psychedelic doses of several compounds, including
psilocybin \cite{schartner2017}. Stoliker et al.\ analysed psilocybin fMRI/EEG
across rest, meditation, music, and movie contexts and argued that variability
can remain structured and context-sensitive rather than randomly
disorganised; stronger subjective effects accompanied more cohesive
context-specific trajectories \cite{stoliker2026}. Bredenberg et al.\ modelled
a shift toward top-down generative influence and showed that structured
internal weighting can produce variability and distortions distinct from
additive noise controls \cite{bredenberg2026}.

None of these sources is a working-memory RNN. Their joint implication is
nonetheless clear: candidate perturbations should include gain-like and
weighting manipulations, and generic noise should be retained as a control
rather than treated as the default psychedelic mechanism.

\subsection{Serotonergic and gain mechanisms relevant to working memory}

Cellular and circuit evidence makes response gain a particularly tractable
bridge. In rat prefrontal pyramidal neurons, serotonin increased the slope of
the firing-rate/current relationship through postsynaptic 5-HT2 mechanisms
\cite{zhang2005gain}. In monkeys performing oculomotor delayed response, local
5-HT2A stimulation and blockade altered delay-period spatial tuning in a
dose-dependent and non-monotonic manner involving both putative excitatory and
inhibitory cells \cite{williams2002ht2a}. Cano-Colino and colleagues used a
receptor-informed spatial working-memory attractor model to show that
serotonergic manipulations can produce distinct failure modes---delay decay,
spurious bumps, distractibility, or perseveration---and that isolated 5-HT2A
agonism increased distractor resistance in their model rather than simple
distractor capture \cite{canocolino2013}. That result is an important
constraint: behavioural distractor vulnerability under psilocybin should not be
equated a priori with a receptor claim that 5-HT2A agonism increases distractor
gain.

Direct psilocin work strengthens heterogeneity. Schmitz et al.\ found that
psilocin excited identified medial-prefrontal 5-HT2A neurons while producing
mixed increases, decreases, and null effects across the broader pyramidal
population \cite{schmitz2025}. Human cortical-propagation analyses further
suggest that psilocybin effects depend on the heterogeneous 5-HT2A receptor
landscape \cite{makimarttunen2026}. These findings motivate testing both global
and unit-wise heterogeneous drive-gain analogues, while treating sensory,
distractor, and recurrent weightings as separate competing translations.

\subsection{Evidence gaps and modelling implications}

Across this core set, no reviewed source both trains a working-memory RNN and
tests which literature-motivated dynamical perturbation recovers selective
acute psilocybin behavioural signatures more convincingly than generic
disruption. To our knowledge, that combined question remains open. The
literature therefore justifies the present programme as computational
sufficiency and comparative mechanism, not as biological identification:
build competent delay-period networks; define selective behavioural contrasts
from human task evidence; compare gain, afferent, recurrent, persistence, and
timescale interventions against a native baseline and, later, against matched
Gaussian disruption; and explain any selective match through population
dynamics rather than accuracy alone.
